\centerline{\bf \Huge Вписанный четырёхугольник}

\begin{flushright}
        \scriptsize{Эпизод 2. Зверь}
\end{flushright}

\begin{figure}[h]
    \centering
    \includegraphics[width=1\linewidth]{Голодные игры 8+/G2.png}
    \caption{Критерий вписанности}
    \label{fig:placeholder}
\end{figure}

\problem{1}
\textbf{Лемма Фусса.} Окружности $\omega_1$ и $\omega_2$ пересекаются в точках $A$ и $B$. На $\omega_1$ выбираются точки $X \neq A$ и $Y \neq B$. Лучи $X A$ и $Y B$ пересекают окружность $\omega_2$ в точках $X_1$ и $Y_1$ соответственно. Докажите, что $X Y \| X_1 Y_1$.

\problem{2}
На сторонах $A B, B C$ и $C A$ треугольника $A B C$ выбраны такие точки $C_1, A_1$ и $B_1$, а внутри него --- точку $P$. Эту точку соединили с точками на сторонах треугольника так, что образовалось три четырехугольника. Известно, что два из этих четырехугольников --- вписанные. Докажите, что тогда и третий четырехугольник также вписанный.

\problem{3}
В треугольнике $A B C$ проведена высота $A H$. Из точки $H$ опущены перпендикуляры $H P$ и $H Q$ на стороны $A B$ и $A C$ соответственно. Докажите, что четырехугольник $B P Q C$ --- вписанный.

\problem{4}
Пусть $I$ --- центр вписанной окружности остроугольного треугольника $A B C, M$ и $N$ --- точки касания вписанной окружности сторон $A B$ и $B C$ соответственно. Через точку $I$ проведена прямая, параллельная стороне $A C$, и на нее опущены перпендикуляры $A P$ и $C Q$. Докажите, что точки $M, N, P$ и $Q$ лежат на одной окружности.

\problem{5}
\textbf{Прямая Симсона.} Дан остроугольный треугольник $A B C$. Точка $P$ выбирается на меньшей дуге $B C$ описанной около него окружности. Докажите, что проекции точки $P$ на стороны треугольника лежат на одной прямой.

\newpage

\hspace{3mm}

\problem{6}
Дан остроугольный треугольник $A B C$. Серединный перпендикуляр к отрезку $A B$ пересекает отрезок $A C$ в точке $X$, а серединный перпендикуляр к отрезку $A C$ пересекает продолжение отрезка $A B$ за точку $B$ в точке $Y$. Докажите, что точки $B, C, X$ и $Y$ лежат на одной окружности.

\problem{7}
Выпуклый четырехугольник $A B C D$ таков, что $A C=D C$ и $\angle D B C+\angle A B C=180^{\circ}$. Докажите, что этот четырехугольник - вписанный.

\problem{8}
Окружности $\omega_1$ и $\omega_2$ пересекаются в точках $A$ и $B$. По окружности $\omega_1$ с одинаковыми угловыми скоростями против часовой стрелки движутся две точки $X$ и $Y$. Прямые $A X$ и $B Y$ повторно пересекают окружность $\omega_2$ в точках $X_1$ и $Y_1$ соответственно. Докажите, что прямая $X_1 Y_1$ касается некоторой фиксированной окружности.